\section{The shorter inherited presentations, with the new tuple}

For users prioritizing a concise system, the seven-equation 59-operation
presentation remains useful. Its sixteen positive variables are
\[
 J,F,\alpha,z,f,h,i,o,j,s,w,\tau,\eta,\zeta,y,\gamma.
\]
Use $q,X,Y,E,k,c,a,H,\Delta,C,R$ as defined in Section~\ref{sec:norm},
but in this section put
\[
 D=X+ac+\gamma H,\quad L_0=EkY,\quad V=of-c,\quad
 U_0=jc-R,\quad I=ic^2,\quad A_0=\Delta(f^2-1).
\]
Then the complete system is
\begin{equation}\label{eq:system59}
 \begin{aligned}
 \tau^2-L_0(L_0+k)&=1, & D^2-\Delta c^2&=1,\\
 A_0(V^2-y^2)+y^2&=1, & k-hE-R&=1,\\
 (K+w)C&=F+z(q-1), & I^2&=A_0,\\
 U_0&=V.
 \end{aligned}
\end{equation}
Here $z$ is the shifted transport coordinate. The largest residual degree
is 24. Graph lifts can lower this degree at the same arithmetic count but
add witnesses and equations; those existing tradeoffs are not recounted as
new results.

For one polynomial at the smallest inherited arithmetic count, retain the
fifteen positive coordinates
\[
 J,F,\alpha,z,f,h,i,T,s,w,\tau,\eta,\zeta,y,\gamma.
\]
Use the same common abbreviations and set
\begin{align*}
 V&=c(Tf-1)-Rf^2,&K_a&=(i\Delta c^2)^2,\\
 N_0&=\tau^2-L_0(L_0+k),&N_1&=D^2-\Delta c^2,\\
 N_a&=K_a(V^2-y^2)+y^2,&N_s&=\Delta f^2-K_a,\\
 N_k&=k-hE-R,&N_t&=(K+w)C+q-F-z(q-1).
\end{align*}
Here $z$ is the unshifted coordinate. The equation is
\begin{equation}\boxed{
 N_0N_1N_aN_kN_tN_s-\Delta=0.}\label{eq:p67}
\end{equation}
It has 67 operations and exact degree 155. Its complete inherited
unit-sign and Pell proof is included in \path{proofs/INHERITED_P67_PROOF.md};
the new local contribution to this instance is the smaller complete
coefficient tuple. We do not claim a 66-operation polynomial.

The strongest caution in this proof is useful more generally: weak
transport initially permits $C=0$ and a negative unit sign. Discarding
that boundary before recovering the main Pell index would invalidate the
argument. The source proofs retain the boundary, recover the target index,
eliminate the negative sign, and only then conclude $C>0$.

\section{Verification, reproducibility, and provenance}

\subsection{What the finite checks establish}

The archive is intended to be an additive, reviewable research package.
It includes the actual four coefficients, source machine transitions,
literal arithmetic circuits, exact coefficient recipes, source pins,
mathematical proofs, standard-library Python programs, LaTeX source, and
the compiled PDF. No network request is required to replay the main checks.

\begin{longtable}{@{}>{\raggedright\arraybackslash}p{0.25\textwidth}>{\raggedright\arraybackslash}p{0.65\textwidth}@{}}
\toprule
Component & Verification and its scope\\
\midrule
\endhead
Finite language & Reconstruct all 943 ordered tiles from the 430 literal
machine actions, and verify the complete tile-list hash.\\
Coefficient adapter & Independently reconstruct all four full integers;
check 66,062 native-basis identities, exact populations and congruences,
carry thresholds, support intersections, primality, orders, lifting, and
actual unit residues.\\
Positive norms & Prove the quadratic-remainder identity algebraically;
check exact eight-conjugate evaluations on signed fixtures, predecessor
agreement, transport substitutions, liveness and operation ledgers, and
full polynomial degree diagnostics in two finite fields.\\
Nine-variable evaluator & Check 675 exact symbolic correction terms,
12,032 intermediate-value comparisons, and 128 affine-coordinate checks;
verify all supplied inputs/gates are live.\\
Jones constants & Verify actual compiler payload hashes and bit lengths,
all 402 strict index bounds, the complete zero-record convention, and
the actual $z_J$ payload.\\
Release & Run the package from its delivered relative layout, retain
receipts and hashes, compile the TeX document, and inspect rendered pages.\\
\bottomrule
\end{longtable}

These checks do not test whether the actual consistency polynomial has a
zero. They also do not replace the infinite-direction compiler or Pell
arguments. Random agreement is corroboration, whereas a polynomial
identity is proved by the displayed algebra or exact sparse comparison.
Finite-field degree checks corroborate the uniform leading-term proofs;
they are not the sole basis for the exact degree claims.

\subsection{Commands and bindings}

From the package root, run
\begin{verbatim}
python3 verify.py
python3 reproduce.py
\end{verbatim}
The first performs the release checks. The second regenerates the new
coefficients and circuits before checking them. Both use Python 3.10 or
later with the standard library. To rebuild the article:
\begin{verbatim}
pdflatex -interaction=nonstopmode -halt-on-error report.tex
pdflatex -interaction=nonstopmode -halt-on-error report.tex
\end{verbatim}

Every recommended circuit identifies its variable domain and coefficient
binding. The names $K_0$ and \texttt{Kconstant} alias $K$; \texttt{Bm1}
aliases $\beta$; \texttt{MF} aliases the shifted $M_F$.
Historical files in \texttt{source/} or \texttt{upstream/} retain their
original coefficients and are explicitly comparisons. The recommended
nine-variable circuit has its own newly derived index constants. No old
Jones index is silently paired with a different quartic.

The only use of the remote repository was read-only source inspection.
No commit, branch, pull request, issue, or other remote object was created.
Source license and attribution files are preserved. A proposed placement
is a new directory under the user's existing
\texttt{con-zfc-diophantine-forms} collection; the user can incorporate it
using the repository's preferred intake process.

\section{What remains worth trying}

Several plausible shortcuts fail to establish what they appear to promise.
Deleting $N_k=1$ and its coordinate $h$ would appear to give eleven positive
variables. But that equation supplies $k\equiv R+1\pmod E$, hence the
first-Pell-index congruence $2n\equiv R+1\pmod E$. After main-index
recovery it forces the exact relation needed by the half-binomial population
argument. The strict ratio alone gives bounds, not this congruence. No
eleven-variable reduction is claimed.

Likewise, replacing $\beta$ alone by a smaller base, replacing a main
congruence-carrying root by its absolute value, or merging radical tests
without checking repeated or negative radicands can introduce extraneous
solutions. Coefficient numerical audits establish an interface only when
accompanied by its completeness controller and semantic proof.

The next substantial gains are likely to require one of three more
structural advances. A smaller quadratic-residual parent would reduce the
$5^{\nu+1}$ exponent in the Jones construction. A new way to encode the
first-index congruence could remove an actual moderate-degree witness.
A different finite computation language or more compact native geometry
could reduce the fixed ports beyond the conservative layout retained here.
These are specific remaining proof obligations, not asserted results.

For incorporation now, the 59-operation system is the shortest transparent
display, the 103-operation twelve-variable system is a useful balance of
coordinates and evaluation size, and the 1582-operation nine-positive-
variable polynomial is the coordinate-minimizing recommendation. Each has
an explicit fixed instance and the requested ZFC-provable equivalence.

\clearpage
\begin{thebibliography}{9}
\small\raggedright

\bibitem{collection}
Vladimir Reshetnikov and contributors.
\emph{Con(ZFC) in Diophantine form}, audited wave-4 collection.
ProveIt, commit \texttt{8c6eec9d19ef80538c8e61ff6e0180040583446a}.
\href{https://github.com/VladimirReshetnikov/ProveIt/tree/8c6eec9d19ef80538c8e61ff6e0180040583446a/SetTheory/Cardinals/docs/reports/hilbert-tenth-problem/con-zfc-diophantine-forms}{Pinned collection}.
Local copy: \path{evidence/COMPARISON_AT_PIN.md}.

\bibitem{norms}
ProveIt report packages 24 and 30.
\emph{Smaller explicit Diophantine forms of ZFC consistency};
\emph{A smaller twelve-positive-variable synthesis}.
Pinned source proofs and polynomial circuits are recorded in
\path{positive/SOURCE_PINS.json}.

\bibitem{coefficients}
ProveIt report packages 24, 25, and 27.
Mixed-radix native compiler, Wang943 language, arbitrary-odd-dimension
prime padding, and complete coefficient adapters.
Exact source paths and SHA-256 values:
\path{coefficients/SOURCES.json}; retained proofs:
\path{coefficients/source_proofs/}.

\bibitem{quartic}
ProveIt report package 30.
\emph{A 31-positive-variable quartic by eliminating a packing register}
and \emph{The compact nine-variable construction}.
Exact input manifest: \path{nine/SOURCE_MANIFEST.json}.

\bibitem{endpoint}
ProveIt. \emph{KEndpoint215.lean}.
\href{https://github.com/VladimirReshetnikov/ProveIt/blob/8c6eec9d19ef80538c8e61ff6e0180040583446a/Computability/BusyBeaver/Lean/BusyBeaver/ZFC/KEndpoint215.lean}{Pinned formal endpoint}.
Inspected, not rebuilt in this work; provenance copy in \path{proofs/}.

\bibitem{jones}
James P. Jones. Universal Diophantine equation.
\emph{Journal of Symbolic Logic} 47(3), 549--571 (1982).
\href{https://doi.org/10.2307/2273588}{doi:10.2307/2273588}.

\bibitem{repojones}
ProveIt. Corrected Jones manuscript and parameterized nine-variable
development. \path{Computability/HilbertTenthProblem/Papers/1982/} and
\path{Computability/HilbertTenthProblem/Lean/Diophantine/Paper1982/}.
The present composition uses the ordinary source theorem as identified in
the retained package-30 proof, not a newly rebuilt formal specialization.

\bibitem{sun}
Zhi-Wei Sun. \emph{Ten unknowns for Hilbert's tenth problem over the integers}.
arXiv:2609.01594v2 (2 September 2026).
\href{https://arxiv.org/html/2609.01594v2}{Versioned primary source}.
This general theorem is discussed for comparison only and is not instantiated
as a ten-integer deliverable here.

\end{thebibliography}
